\textit{SINDy and relatives.} SINDy~\cite{brunton2015discovering} and its PDE counterpart PDE-FIND~\cite{rudy2016data} use sparse regression over a candidate library. SINDy-PI~\cite{kaheman2020sindy} targets implicit dynamics with a parallel formulation reported to be robust to noise, and Ensemble-SINDy~\cite{fasel2021ensemble} bags models for the low-data, high-noise regime; neither is evaluated here. PySINDy~\cite{silva2020pysindy} packages these methods with several differentiation options.

\textit{Numerical differentiation of noisy data.} Savitzky--Golay local polynomial fitting~\cite{savitzky1964smoothing} and TV-regularised differentiation~\cite{chartrand2011numerical}, which allows non-smooth derivatives, are the classical smoothing differentiators. van Breugel et al.~\cite{breugel2020numerical} cast the choice of method and parameters as a multi-objective optimisation and stress that parameter selection is the practical difficulty.

\textit{Weak and integral formulations.} Schaeffer and McCalla~\cite{schaeffer2017sparse} identify dynamics from integral terms, Reinbold et al.~\cite{reinbold2020using} use a weak form to discover models from noisy and incomplete spatiotemporal data, and Weak SINDy integrates the dynamics against test functions so that pointwise derivative approximations are replaced by linear transformations of the data, for ordinary~\cite{messenger2021weak} and partial~\cite{messenger2020weak} differential equations. The ODE paper~\cite{messenger2021weak} is the direct precedent of this study: it reports that the weak form improves on standard SINDy by orders of magnitude and identifies the correct terms in both small- and large-noise regimes.

\textit{Gap.} Weak SINDy~\cite{messenger2021weak} is compared with standard SINDy in its own paper, and the differentiation literature~\cite{breugel2020numerical,chartrand2011numerical} evaluates derivative estimates, not the equations recovered from them. We did not find a head-to-head of all five derivative sources with a shared STLSQ back end, one tuning protocol and the exact-support rate as the headline metric; this study provides one at small scale. Our weak form is far simpler than WSINDy (one test-function family, fixed width, no reweighting), so the results do not bear on that algorithm as published.
